\section{Objective and the precise outcome}
The Turn 5 handoff asks for positivity of the complete original first-half
pair sum at a prescribed prime, not for another family constructed after useful
factors are supplied. We make that attempt here. The exact pair source and its
$2\tau_{\rm div}(h)$ fibres are inherited from \cite[\S\S3--4]{Turn5}; the
Type I/II coordinates are classical \cite[\S2]{ET}. They are rederived below so
that the new estimates have a complete arithmetic return.

The principal-character isolation attempt fails quantitatively: for every prime
$p\equiv1\pmod4$ with $p\ge2^{20}$, its complete unweighted lower estimate is at
most $-p\log p/24$. More generally this remains true when at most $J$ even characters are retained
at each shell, once $p\ge2^{20}J^3$, including arbitrary shellwise subgroup
observations of index at most $J$. An all-exponent refinement extends this
obstruction to every fixed sublinear mode budget, with an explicit onset.
This is an obstruction to that specified estimate, not to ES or to every possible
signed estimate. An original finite prime also defeats \emph{every} nonnegative
shell reweighting of the principal-only estimate.

The surviving calculation retains nonprincipal even-character information
through unsaturated cosets and uses an exact integer collision bound. It returns
actual states at the four bounded principal-only failures, including a cubic
quotient whose two occupied coarse coefficients are both eight. No theorem
that these improved bounds are positive at some shell of every prime is proved.
No historical-priority claim is made for the estimates or the classical
Cauchy--Schwarz, balancing, Fourier, and divisor arguments they use.

\section{Complete source, first-half completeness, and original returns}
Let $p\equiv1\pmod4$ be prime and put
\[
 I_p=\{a\in\mathbb Z:p/4<a<p/2\},\qquad R_a=4a-p.
\]
Thus $R_a=3,7,\ldots,p-2$, and $\gcd(pa,R_a)=1$.
Write $a=\prod_\ell\ell^{e_\ell}$. The complete original divisor source is
\[
 \beta_\ell\in[-e_\ell,e_\ell]_{\mathbb Z},\qquad
 u=\prod_\ell\ell^{e_\ell+\beta_\ell}\mid a^2.            \tag{1}
\]
Its inverse is $\beta_\ell=v_\ell(u)-e_\ell$. The gates and normalization are
\[
 E:R_a\mid4u+1,\qquad M:R_a\mid4u+p,
\]
\[
 d_0=\gcd(a,u),\quad (h,r,s)=(d_0^2/u,u/d_0,a/d_0),
 \quad a=hrs,\quad u=hr^2,\quad\gcd(r,s)=1.
\]
Prime valuations prove the asserted positive integers. The ordered returns are
\[
 E:\ \kappa=(pr+s)/R_a,\quad(x,y,z)=(a,hs\kappa,phr\kappa),
                                                               \tag{2}
\]
\[
 M:\ \lambda=(r+s)/R_a,\quad(x,y,z)=(a,phs\lambda,phr\lambda).
                                                               \tag{3}
\]
The gates and the fact that $h,r,s$ are units modulo $R_a$ prove the integer
quotients. Clearing denominators proves $4/p=1/x+1/y+1/z$. The ordered divisor
inverses are respectively
\[
 u=\frac{a^2}{R_ay-pa},\qquad
 u=\frac{pa^2}{R_ay-pa}.                                      \tag{4}
\]
The middle involution $u\mapsto a^2/u$ swaps the last two denominators;
exterior orientation retains its unique $p$-divisible denominator last.

For completeness, every positive solution has a distinguished first-half
coordinate. Order $x\le y\le z$. Then $p/4<x\le3p/4<p$, and at least one,
but not all, denominators are multiples of $p$. In Type II, $y,z\ge p$ force
$x<p/2$. In Type I write the two non-$p$ denominators as $dr,ds$, with
$r\le s$ and $\gcd(r,s)=1$, and the third as $pk$. The equation gives
$p\mid4k-1$; put $t=(4k-1)/p\equiv3\pmod4$. Then
\[
 td rs=k(r+s),\quad k=h rs,\quad r+s=t\lambda,\quad d=h\lambda.
\]
Here $rs\mid k$ and $\gcd(t,h)=1$ justify the two divisions.
The case $r=s$ would give $t\mid2$ and is impossible. For $r<s$,
\[
 p-2hr\lambda
 =\frac{2hr(s-r)-1}{t}>0.
\]
So $hr\lambda$ is a first-half denominator. Finally
\[
 (R_ay-pa)(R_az-pa)=p^2a^2                                 \tag{5}
\]
classifies a solution at this denominator. Both factors are positive; their
$p$-valuations are $(0,2),(1,1),(2,0)$. The required congruence modulo $R_a$
and the divisor relation recover exactly (1)--(4), with the recorded exterior
orientation or either middle order. These are the inherited Type I/II
coordinates \cite[Propositions 2.2 and 2.6]{ET}, not a new parametrization.

\subsection{The common pair source}
Retain actual divisors of $pa$ as $d=p^\epsilon v$, where $v\mid a$ and
$\epsilon\in\{0,1\}$. Put
\[
 \mathcal P_a=\{(b,c):b,c\mid pa,\ R_a\mid b+c\},\qquad T_a=|\mathcal P_a|.
                                                               \tag{6}
\]
A primitive labelled M pair $(r,s)$ has the full fibre
\[
 (b,c)=(p^\epsilon d r,p^\epsilon d s),\quad d\mid h,
 \quad \epsilon\in\{0,1\};                                \tag{7}
\]
an E pair has fibre
\[
 (pdr,ds)\ \hbox{or}\ (ds,pdr),\quad d\mid h.              \tag{8}
\]
To invert, divide $(b,c)$ by its gcd. If neither coprime part contains $p$, use
M and retain the gcd's $p$-bit. Otherwise remove its one $p$, designate that side
$r$, and use E, retaining the side. The lcm condition gives $drs\mid a$ and
$h=a/(rs)$. Thus every original state has exactly $2\tau_{\rm div}(h)$ preimages.
At fixed $a$ the map $(r,s,h)\mapsto u=hr^2$ is the original divisor inverse.
Consequently, for $F_a=|E_a|+|M_a|$,
\[
 2F_a\le T_a\le2\tau_{\rm div}(a)F_a,\qquad
 \operatorname{ES}(p)\ \Longleftrightarrow\ \sum_{a\in I_p}T_a>0. \tag{9}
\]
The full exact-state M\"obius inversion remains in \cite[\S4, ``Complete primitive count'']{Turn5}; it
is not an additional prerequisite for using (9). Here $\tau_{\rm div}$ is the
divisor-count function, not the absent-support element $\tau$.

\begin{lemma}[Original fourfold orbits]\label{lem:four}
For every actual shell, $T_a\in4\mathbb Z_{\ge0}$. The involutions
\[
 S(b,c)=(c,b),\qquad I(b,c)=(pa/b,pa/c)                    \tag{10}
\]
generate a free Klein four action on $\mathcal P_a$.
\end{lemma}
\begin{proof}
They commute and preserve (6). A fixed point of $S$ would give $R_a\mid2b$,
contrary to coprimality and $R_a\ge3$. A fixed point of $I$ would give $b^2=pa$,
contrary to $v_p(pa)=1$. A fixed point of $SI$ gives $bc=pa$ and $b\equiv-c$.
Since $pa\equiv4a^2\pmod{R_a}$, it would make $-1$ a square modulo $R_a$.
This is impossible when $R_a\equiv3\pmod4$: a prime $3\pmod4$ occurs to odd
exponent in $R_a$ and $-1$ is a nonsquare there. Thus all orbits have four members.
In (7), $S$ swaps $r,s$ and $I$ additionally sends $(d,\epsilon)$ to
$(h/d,1-\epsilon)$. In (8), $S$ changes the side and $I$ also replaces $d$ by
$h/d$. These are actions on the original marked fibres, not extra ES states.
\end{proof}

\section{The full-shell principal estimate is uniformly negative}
In $G_a=(\mathbb Z/R_a\mathbb Z)^\times$, let
\[
 c_a(g)=\#\{d\mid pa:d\equiv g\pmod{R_a}\},\quad
 m_a=2\tau_{\rm div}(a),\quad \mathcal E_a=\sum_g c_a(g)^2.
\]
Character orthogonality gives
\[
 T_a=\frac1{\varphi(R_a)}\sum_\chi\chi(-1)
             \left|\sum_{d\mid pa}\chi(d)\right|^2,
 \qquad
 \mathcal E_a=\frac1{\varphi(R_a)}\sum_\chi
             \left|\sum_{d\mid pa}\chi(d)\right|^2.       \tag{11}
\]
All characters, including imprimitive ones, are retained. The first identity is
already present in the historical pair-source transcript \cite{Boolean}; its
use here does not claim a new Fourier reformulation. The principal contribution
is $P_a=m_a^2/\varphi(R_a)$. Giving every nonprincipal term the adverse sign yields
\[
 T_a\ge L_a:=2P_a-\mathcal E_a
 =\frac{8\tau_{\rm div}(a)^2}{\varphi(R_a)}-\mathcal E_a. \tag{12}
\]
This is a valid lower bound, but the next theorem shows why summing it cannot
close the conjecture.

\begin{lemma}[Two elementary explicit bounds]\label{lem:elementary}
For $n\ge1$, $\tau_{\rm div}(n)<4n^{1/3}$. For $x\ge3$,
\[
 \sum_{\substack{r\le x\\r\equiv3\ (4)}}\frac1{\varphi(r)}
 <\frac75\left(1+\frac14\log x\right).                  \tag{13}
\]
\end{lemma}
\begin{proof}
For primes at least eleven, $e+1\le2^e\le q^{e/3}$. For $q=2,3,5,7$, successive
ratios show that $(e+1)/q^{e/3}$ is maximized at $e=3,2,1,1$, respectively.
The four maxima are bounded by $2,3/2,6/5,11/10$; their product is $3.96<4$.
The ratio check after a maximum is $(e+2)^3<q(e+1)^3$ and becomes stronger as
$e$ increases; before the maximum the reverse inequalities hold. This proves
all exponent cases, not only the finitely evaluated maxima.

For the second assertion use the classical divisor identity, also recorded in
\cite[Lemmas 8--10, pp. 11--12]{Farah}. Only this identity and its elementary
expansion are used; no probability or independence assertion from that paper
enters the present proof. Explicitly,
\[
 \frac r{\varphi(r)}=\sum_{d\mid r}\frac{\mu(d)^2}{\varphi(d)}.
\]
Writing $r=dm$, both variables are odd and $m$ is in one specified reduced class
modulo four. For either class,
$\sum_{m\le X}1/m\le1+\frac14\log X$: isolate its first term and bound later
terms by the preceding length-four integrals. Consequently the left side of
(13) is at most
\[
 \left(1+\frac14\log x\right)
 \sum_{d\text{ odd}}\frac{\mu(d)^2}{d\varphi(d)}
 =\left(1+\frac14\log x\right)
 \prod_{q\text{ odd prime}}\left(1+\frac1{q(q-1)}\right).
\]
Absolute convergence justifies this expansion. Its product is at most
\[
 \exp\left(\sum_{j\ge1}\frac1{2j(2j+1)}\right)
 =\exp(1-\log2)=\frac e2<\frac75.
\]
The last strict inequality follows from the factorial series for $e<14/5$.
\end{proof}

\begin{theorem}[A complete pointwise failure of principal domination]
\label{thm:negative}
For every prime $p\equiv1\pmod4$ with $p\ge2^{20}$,
\[
 \boxed{\sum_{a\in I_p}L_a\le-\frac1{24}p\log p.}        \tag{14}
\]
\end{theorem}
\begin{proof}
Put $\mathcal P=\sum_aP_a$ and $\mathcal D=\sum_a\mathcal E_a$.
Because $a<p/2$ and $2^{-2/3}<2/3$, Lemma~\ref{lem:elementary} gives
\[
 \mathcal P<\frac{224}{15}p^{2/3}(\log p+4).             \tag{15}
\]
On the other hand, $\mathcal E_a\ge m_a=2\tau_{\rm div}(a)$ from the original
diagonal pairs. Let $Y=\lfloor\sqrt p/2\rfloor$. Every divisor $d\le Y$ of
$a$ has a distinct complementary divisor $a/d>Y$, since $a>p/4\ge Y^2$.
The interval $I_p$ contains more than $p/(4d)-1$ multiples of $d$, so
\[
 \mathcal D\ge2\sum_a\tau_{\rm div}(a)
 >p\sum_{d=1}^Y\frac1d-4Y
 >p(\tfrac12\log p-\log2)-2\sqrt p.                    \tag{16}
\]
The harmonic bound is $\sum_{d\le Y}1/d\ge\log(Y+1)$, and $Y+1>\sqrt p/2$.

For $p\ge2^{20}$, one has $p^{1/3}>100$, $\log p>40/3$, $\log2<3/4$ and
$2/\sqrt p\le1/512$. Equations (15)--(16) imply
\[
 \frac{2\mathcal P}{p\log p}
 <\frac{448}{15}\frac1{100}\frac{13}{10}
 =\frac{728}{1875},\qquad
 \frac{\mathcal D}{p\log p}>\frac7{16}.
\]
Finally $7/16-728/1875=1477/30000>1/24$. This proves (14). The elementary
logarithm bounds follow either from integration or the rational series retained
in the certificate. No distribution-of-primes theorem or ES occupancy enters.
\end{proof}

\begin{corollary}[Adaptive finite quotient resolution]\label{cor:index}
For each shell choose any subgroup $B_a\le G_a$ containing $-1$, with
$[G_a:B_a]\le J$. Let $m_{a,C}=\sum_{g\in C}c_a(g)$ and put
\[
 L_{a,B_a}=\frac2{|B_a|}\sum_{C\in G_a/B_a}m_{a,C}^2-\mathcal E_a.
                                                               \tag{17}
\]
For every integer $J\ge1$ and prime $p\equiv1\pmod4$ with $p\ge2^{20}J^3$,
\[
 \sum_a L_{a,B_a}\le-\frac1{24}p\log p.                 \tag{18}
\]
The subgroups may depend on the actual factorization of each shell.
More generally, let $\Sigma_a$ be any at most $J$ even characters, with conjugate
pairs retained when a real observation is intended, and put
\[
 L_{a,\Sigma_a}=\frac2{\varphi(R_a)}
 \sum_{\chi\in\Sigma_a}\left|\sum_{d\mid pa}\chi(d)\right|^2-\mathcal E_a.
                                                               \tag{18a}
\]
Then $T_a\ge L_{a,\Sigma_a}$ and the same summed upper bound (18) holds.
\end{corollary}
\begin{proof}
Let $\mathcal E_+$ be the full even-character energy. Equation (11) gives
$T_a=2\mathcal E_+-\mathcal E_a$. Retaining any part of that nonnegative energy
proves (18a). Each squared character sum is at most $m_a^2$, so
$L_{a,\Sigma_a}\le2JP_a-\mathcal E_a$. The previous proof applies because
$p^{1/3}>100J$. The subgroup version is $\Sigma_a=B_a^\perp$, whose cardinality
is its index. Equivalently use $\sum_Cm_{a,C}^2\le m_a^2$ in (17).
\end{proof}
The same argument permits positive weights $w_0\le w_a\le w_1$ after replacing
$J$ by $Jw_1/w_0$ in the height condition, and yields the upper bound
$-w_0p\log p/24$. It does not cover truncating all negative bounds to zero,
weights without a positive lower bound, unbounded-index observations, or the
integer refinement below. Those restrictions are part of the theorem.

\subsection{An effective obstruction at every fixed sublinear spectral scale}
The power $J^3$ above is not an intrinsic barrier. For an integer $m\ge3$, put
\[
 K_m=\prod_{\substack{q<2^m\\q\ \mathrm{prime}}}
       \max_{0\le e<2m}\frac{(e+1)^m}{q^e},\qquad
 C_m=\min\{C\in\mathbb Z_{>0}:C^m\ge K_m\}.             \tag{18b}
\]
These are finite, exactly computable constants. For every integer $a\ge1$,
\[
 \tau_{\rm div}(a)\le C_ma^{1/m}.
\]
For $q\ge2^m$, use $e+1\le2^e$. For a smaller prime, the consecutive ratio
is decreasing, and once $e+1\ge2m$,
\[
 \left(1+\frac1{e+1}\right)^m
 \le\left(1+\frac1{2m}\right)^m
 \le\sum_{j=0}^m\frac1{2^j j!}<2\le q.
\]
The maximum is therefore among the stated $2m$ exponents. Multiplication proves
the all-integer bound.

\begin{theorem}[Effective sublinear-resolution obstruction]\label{thm:sublinear}
Retain any at most $J$ even characters at each original shell. For every prime
$p\equiv1\pmod4$ such that
\[
 p\ge2^{20},\qquad p^{1-2/m}\ge12JC_m^2,
\]
\[
 \boxed{\sum_{a\in I_p}L_{a,\Sigma_a}\le-\frac18p\log p.} \tag{18c}
\]
For $m=4$, $C_4=9$, so it suffices that
\[
 \boxed{p\ge\max\{2^{20},944784J^2\}.}                 \tag{18d}
\]
\end{theorem}
\begin{proof}
The reciprocal-totient estimate (13) gives
\[
 \mathcal P\le\frac75 C_m^2\,2^{-2/m}p^{2/m}(\log p+4).
\]
As before $\mathcal D/(p\log p)>7/16$ and
$1+4/\log p<13/10$. The capacity inequality therefore gives
\[
 \frac{2J\mathcal P}{p\log p}<\frac{91}{300},\qquad
 \frac7{16}-\frac{91}{300}=\frac{161}{1200}>\frac18.
\]
This proves (18c). At $m=4$, the primes $2,3,5,7,11,13$ have maximizing
exponents $5,3,2,1,1,1$, and
\[
 K_4=\frac{127401984}{25025},\qquad8^4<K_4<9^4.
\]
Thus $C_4=9$, and $12C_4^2=972$, whose square is $944784$.
\end{proof}
For any fixed $\delta>0$, choose $m$ with $2/m<\delta$. The theorem then
excludes a positive raw sum from any per-shell budget at most $p^{1-\delta}$
for every
\[
 p\ge\max\{2^{20},(12C_m^2)^{1/(\delta-2/m)}\}.
\]
This is a bound on the stated adverse-sign estimate, not a lower bound on the
number of modes needed to locate an ES witness by any method. Zero truncation,
integer balancing and arithmetic control of the omitted signs remain outside
this obstruction.

The diagonal mass in (16) is exactly cancelled in the true signed formula,
since $\sum_\chi\chi(-1)=0$. In particular one may write (11) with each squared
character sum replaced by $|\sum_{d\mid pa}\chi(d)|^2-m_a$. That exact removal
does not itself bound the remaining signed terms. The theorem rules out the
specific adverse-sign replacement, not arithmetic cancellation in the original sum.

\section{Retaining nonprincipal information on unsaturated cosets}
Fix an actual shell and a chosen subgroup $B\le G_a$ with $-1\in B$; it is
\emph{not} assumed to be the actual stabilizer of the divisor support. Write
$|B|=2s$. Every coset is partitioned into $s$ antipodal pairs $\{g,-g\}$.
For integer $m\ge0$, write $m=sq+r$ with $0\le r<s$ and define
\[
 \mathfrak b_s(m)=sq^2+r(2q+1).                            \tag{19}
\]
\begin{theorem}[Integer collision lower bound]\label{thm:integer}
With $m_C=\sum_{g\in C}c_a(g)$,
\[
 \boxed{T_a\ge\mathcal L_{a,B}:=
 \sum_{C\in G_a/B}\mathfrak b_s(m_C)-\mathcal E_a
 \ge L_{a,B}.}                                           \tag{20}
\]
Equality in the first inequality holds exactly when the antipodal-pair totals
within every $B$-coset differ by at most one. Moreover
\[
 \boxed{T_a\ge4\max\{0,\lceil\mathcal L_{a,B}/4\rceil\}.} \tag{21}
\]
\end{theorem}
\begin{proof}
Put $z_{C,j}=c_a(g_j)+c_a(-g_j)$ for the $s$ pairs in a coset. Expansion gives
\[
 T_a+\mathcal E_a=\sum_C\sum_{j=1}^s z_{C,j}^2.
\]
For nonnegative integers of sum $m_C$, transferring one unit from a coordinate
at least two larger than another strictly decreases the sum of squares.
It therefore has minimum (19), with exactly the stated equality locus.
Real Cauchy--Schwarz gives $\mathfrak b_s(m_C)\ge m_C^2/s$, proving the second
inequality. Lemma~\ref{lem:four} proves (21).
\end{proof}

The retained spectral interpretation is exact. Let $Q_Bc$ be the vector constant
at $m_C/|B|$ on each coset, and $J_-c(g)=c(-g)$. In the explicitly specified
Euclidean residue-coordinate metric, $Q_B$ is an orthogonal projection and
$J_-Q_B=Q_B$. Therefore
\[
 T_a=\|Q_Bc_a\|^2+
 \langle c_a-Q_Bc_a,J_-(c_a-Q_Bc_a)\rangle,
\]
which yields (17) without assuming that the support is saturated. Fourier
expansion of $Q_B$ retains \emph{all} characters trivial on $B$; they are even.
For $B=G_a$ this is the principal estimate. For $B=\{1,-1\}$ it is the exact
answer, so that terminal choice is not presented as a reduced forcing theorem.

The bound uses coset counts and equal-residue collisions, not negative-target
membership. When it is positive, enumerate the original residue buckets until
one bucket and its negative are both nonempty, then apply (7)--(8) and (2)--(4).
There are exactly $2\tau_{\rm div}(a)$ original words, so this finite return
terminates under the proved positive bound. This is a terminating return from a
\emph{verified inequality}, not a proof that the inequality always holds.

\subsection{Original fibres and metrics}
The map from the free integer module on original divisors to residue coefficients
sums inside each actual residue fibre. Its kernel is generated by
$e_d-e_{d_*}$, with $d_*$ the least original divisor in that fibre. Retaining the
sum and all nonselected coordinates reconstructs the selected coordinate by
subtraction. The same construction applies to the residue-to-coset map.
Thus neither projection is claimed injective. In the word-counting metric, the
minimum-norm section of residue data $y_g$ assigns $y_g/c_a(g)$ to each word in
that nonempty fibre and has squared norm $\sum y_g^2/c_a(g)$. This is not the
unweighted residue metric used to evaluate the explicitly defined collision
energy. In particular coset summation is not declared isometric.

A positive integer coarse coefficient that reduces to zero modulo two remains
supported with its original fibre. An empty fibre is absent. The proof of (20)
uses the integer counts before such reduction. It never infers a fine target
from coarse membership alone; the retained collision term supplies the bound.

\subsection{The complement-aware refinement}
There is further original information which can be retained without inspecting
the target. Put $P=pa\bmod R_a=(2a)^2\bmod R_a$. Divisor complementation gives
\[
 c_a(g)=c_a(P/g).                                         \tag{22}
\]
If $g^2=P$, the coefficient $c_a(g)$ is even: the fixed residue fibre is paired
by $d\mapsto pa/d$, which has no fixed original word.

Complementation pairs the $B$-cosets. A distinct pair $C,C'$ has equal masses
$m$ and contributes at least $2\mathfrak b_s(m)$ to $T_a+\mathcal E_a$.
For a self-complementary coset, let $f$ be half the number of its roots of $g^2=P$.
The induced involution on its antipodal pairs has $f$ fixed pairs and
$r=(s-f)/2$ two-pair orbits. There are no extra fixed pairs from $g^2=-P$, since
$-1$ is a nonsquare. The coset mass is even, say $2M$. Its exact lower
contribution is
\[
 \mathfrak c_{r,f}(M)=
 \min_{\substack{z_i,w_j\ge0\text{ integers}\\\sum z_i+\sum w_j=M}}
 \left(2\sum_{i=1}^r z_i^2+4\sum_{j=1}^f w_j^2\right).   \tag{23}
\]
This minimum is executable without an oracle: select the $M$ smallest marginal
costs in the $r$ lists $2,6,10,\ldots$ and $f$ lists $4,12,20,\ldots$, retaining
the list labels. Every selected prefix is feasible. Exchange of a selected larger
cost with an unselected smaller cost proves optimality. A separate dynamic
program checks the finite examples. Summing these orbit bounds and subtracting
$\mathcal E_a$ strengthens (20). It is an actual source-complement refinement,
not an assumption that a general coefficient array satisfies (22).

\section{Finite counterexample to all principal-only shell reweighting}
At the deterministically verified hard prime
\[
 p=87481\equiv121\pmod{840},
\]
complete enumeration of all $21870$ first-half shells proves
\[
 \boxed{L_a\le0\quad\text{for every }a\in I_p.}           \tag{24}
\]
The integer refinement (20) with $B=G_a$ is also nonpositive at every shell.
There are nevertheless 34 occupied shells. Equation (24) proves that no choice
of nonnegative shell weights, even chosen after factorization, makes
$\sum_aw_aL_a$ positive at this prime. This is a counterexample to that proposed
certificate family, not to ES.

At $a=21878=2\cdot10939$, $R=31$, take
\[
 B=\langle-1,2\rangle=\{1,2,4,8,15,16,23,27,29,30\}.
\]
The original mass is eight, the energy eight, and the coarse masses are $(8,0,0)$.
Thus $L_{a,B}=24/5$, $\mathcal L_{a,B}=6$, and (21) forces $T_a\ge8$, attained.
The pair $(1,p)$ returns the original E state
\[
 (h,r,s,\kappa)=(21878,1,1,2822),
\]
\[
 \frac4{87481}=\frac1{21878}+\frac1{61739716}+\frac1{5401052095396}.
\]
The shell is already an elementary endpoint success. Its role here is to refute
the claim that the full principal estimate is necessarily a sufficient detector,
not to claim a previously uncovered prime.

\section{A cubic quotient gives a sharp original lower bound}
At the hard prime $p=944329\equiv169\pmod{840}$, every principal-only $L_a$ is
again nonpositive. Choose the original shell
\[
 a=236098=2\cdot97\cdot1217,\qquad R=63,
\]
and the chosen observation subgroup
\[
 B=\langle-1,13\rangle
 =\{1,8,13,20,22,29,34,41,43, 50,55,62\}.
\]
This has index three; it is not a saturated support or assumed stabilizer.
The original 16 divisors of $pa$ have distinct fine residues, so
$\mathcal E_a=16$. The three coset counts are $(8,8,0)$. Since $s=6$,
\[
 \mathfrak b_6(8)=12,\qquad
 \boxed{\mathcal L_{a,B}=12+12-16=8=T_a.}                  \tag{25}
\]
In comparison, the principal-only lower bound is $-16/9$, and the real
coset bound is $16/3$. If $\omega^2+\omega+1=0$, the two retained nonprincipal
cubic coefficients are $8(1+\omega)$ and its conjugate, each of squared
absolute value 64. Their combined positive energy is $32/9$ after division by
$\varphi(63)$. Twice that energy is exactly the improvement from $-16/9$ to
$16/3$. The integer balancing supplies the remaining sharpening to eight.
Both occupied coarse counts reduce to zero modulo two; their original words
and the integer argument do not disappear.

The source pair $(1,1149248393)=(1,944329\cdot1217)$ returns
\[
 (u,h,r,s,\kappa)=(287331266,194,1217,1,18242038),
\]
\[
 \frac4{944329}=\frac1{236098}+\frac1{3538955372}
               +\frac1{4067138774169717196}.
\]
At the same prime, the occupied shell $a=236094$, $R=47$ has $T_a=60$ and
principal lower bound $-732/23$. Its unit group has order 46. There is no proper
subgroup containing $-1$ with order greater than two. Thus this shell has no
nonterminal subgroup observation of the present kind. Taking the order-two
subgroup computes $T_a$ exactly and is not a new reduced certificate. This
negative control prevents universal success of the quotient hierarchy from
being asserted without proof.

\subsection{A proper selected spectrum resolves the prime-order quotient}
The absence of an intermediate subgroup at $R=47$ does not preclude a smaller
spectral observation. Keep the same original shell
\[
 p=944329,\qquad a=236094=2\cdot3\cdot19^2\cdot109,\qquad R=47.
\]
The element $5$ generates its unit group of order 46. Let
$\zeta=e^{2\pi i/23}$ and $\chi(5)=\zeta$, so $\chi(-1)=1$.
The original 48 divisors give, in these even-character coordinates,
\[
 C(X)=2+X+X^{11}+2X^{12}+3X^{13}+4X^{14}+4X^{15}
       +5X^{16}+5X^{17}+5X^{18}+5X^{19}+4X^{20}+4X^{21}+3X^{22}.
\]
For the original norm polynomial $Q(X)=C(X)C(X^{-1})\bmod(X^{23}-1)$,
\[
 Q_0=192,\qquad
 (Q_1,\ldots,Q_{11})=(187,176,159,138,116,92,70, 50,33,21,14),
 \quad Q_{23-j}=Q_j.
\]
Exact rational intervals give
\[
 \boxed{1009<Q(\zeta)<1010,\qquad24<Q(\zeta^2)<25.}       \tag{25a}
\]
For completeness the interval certificate uses
$\pi=16\arctan(1/5)-4\arctan(1/239)$. The rational tangent formula makes
$\tan(4\arctan(1/5)-\arctan(1/239))=1$, and that angle is in $(0,1)$;
$\pi/2>1$ follows from $\arctan1>1/2$, so the identity has the stated branch.
Use 18 alternating arctangent terms with the next-term remainder. For
$0\le x\le\pi$, use 18 cosine terms and the next-term absolute remainder;
all tail terms are decreasing. Uncertainty in $\pi$ contributes at most the
argument uncertainty, since $|\cos x-\cos y|\le|x-y|$. Evaluating
$Q_0+2\sum_{j=1}^{11}Q_j\cos(2\pi mj/23)$ with reflection into $[0,\pi]$
proves (25a) by rational arithmetic. A second checker uses
$\pi=4(\arctan(1/2)+\arctan(1/3))$, 50 arctangent terms and 21 cosine terms,
and obtains the same strict integer enclosures. No floating-point value is a
proof input. Both polynomial and intervals remain in the JSON certificate.

The collision energy is 132 and the principal squared sum is 2304. Retaining only
$\{1,\chi,\bar\chi\}$ in (18a) gives
\[
 T_a\ge\frac{2Q(\zeta)-732}{23}>\frac{1286}{23}>52,
 \quad\text{hence }T_a\ge56.
\]
Adding $\chi^2,\bar\chi^2$ gives
\[
 \boxed{T_a>\frac{2(1009+24)-732}{23}=58,
         \quad T_a\ge60,}                               \tag{25b}
\]
using the original fourfold action. Direct enumeration, performed separately,
gives $T_a=60$. Thus five of the 23 even-character modes certify the full attained
lower bound. Unlike retaining every even mode, this is a proper observation.
The original pair $(2,327)$ returns the M state
\[
 (u,h,r,s,\lambda)=(1444,361,2,327,7),
\]
\[
 \frac4{944329}=\frac1{236094}+\frac1{780326438241}
                         +\frac1{4772638766}.
\]
The norm polynomial is calculated from the original words; its identity
$Q_0=\mathcal E_a+T_a$ is a verification check, not an assumed target input to
the interval inequality. The retained observation consists of the selected
Fourier values and collision energy; the remaining Fourier coordinates form
its explicit kernel complement. Orthogonal Fourier projection onto a
conjugate-stable set of even modes commutes with $J_-$, which is the identity
on that image. Keeping all omitted coordinates gives its inverse. No coset
saturation, change of coefficient prime, or positivity of the omitted signed
part is postulated.

\section{Finite comparison and the remaining universal inequality}
The complete declared hard universe through $2{,}000{,}000$ has 4519 primes.
The principal-only condition $L_a>0$ at some original shell holds at 4515.
Its exact misses are
\[
 87481,\quad196561,\quad944329,\quad1915201.              \tag{26}
\]
Completeness of the finite test does not depend on an arbitrary residual cutoff:
$\mathcal E_a\ge2\tau_{\rm div}(a)$ makes $L_a>0$ require
$\varphi(R)<4\tau_{\rm div}(a)$. The exact maximum of $\tau_{\rm div}(a)$ over
$1\le a\le999999$ is 240. All 311 possible residuals with $R\equiv3\pmod4$,
$R<2{,}000{,}000$ and $\varphi(R)<960$ are tested; their largest is 1995.
A separate implementation additionally enumerates \emph{every} first-half shell
of all four exceptions, not only this reduced candidate set.

Chosen quotient certificates for the four primes have
$(a,R,|B|,\mathcal L_{a,B},T_a)$ respectively
\[
 (21878,31,10,6,8),\quad(49159,75,20,4,16),
\]
\[
 (236098,63,12,8,8),\quad(478836,143,4,86,144).
\]
The last complement-aware lower bound is 88. At the second prime, even the
integer principal observation at $a=49147$, $R=27$ gives lower bound two where
the real bound is zero. The certificates retain this simpler additional repair.
The basic integer scalar test therefore covers 4516 of the 4519 primes; the
displayed nonprincipal certificates cover the remaining three. These are
comparisons of sufficient tests, not newly solved or previously unverified primes.

The complete small replay checks all 73798 first-half shells at the 211 primes
$p\equiv1\pmod4$, $p\le3000$. It retains 1877004 divisor vectors, 5700 E states,
5268 oriented M states, and 70596 target pairs. The independent implementation
uses literal pair comparisons and trial enumeration of square divisors on that
whole range. Finite computations do not establish a positive lower bound at all
primes. Theorem~\ref{thm:negative} and Corollary~\ref{cor:index} have general
written proofs and do not extrapolate from this scan.

The exact surviving sufficient quantity is
\[
 \sum_{a\in I_p}4\max\{0,\lceil\mathcal L_{a,B_a}/4\rceil\},
                                                               \tag{27}
\]
or its complement-aware refinement. It is bounded above by $\sum_aT_a$ and
returns original witnesses when positive. No proof that (27) is positive for
some nonterminal choices at every prescribed prime is supplied. Likewise no
universal upper bound on the remaining signed Fourier contribution has been
proved. The source's requested universal closure is therefore unresolved.

The next estimate must force enough of the actual even-character
energy without assuming a favourable spectrum, or exploit original off-diagonal correlations before absolute values, or
supply a different arithmetic witness. The preceding factor-supply character
is annihilated in the combined observation \cite[\S4, ``The source character is annihilated'']{Turn5}; repeating
that fact cannot provide the missing positive term. Averaged ES counts
\cite[Theorem 1.1 and Remark 1.2]{ET} are not converted into a pointwise lower
bound. There is no Lean build, independent mathematical review, or historical
priority assertion in this tranche.
